\section{Physical equivariance and the complete boundary}\label{sec:boundary56}
Return now to a compact Hausdorff group $H$ with a finite dense generating alphabet. The output sets may still be norm-compact convex subsets of real Banach spaces. Purification removes transition randomness; a further averaging removes hidden relations that have no physical counterpart.

\begin{theorem}[Physical-equivariant purification]\label{thm:equivariant56}
Every stationary $k$-label realization of error at most $\varepsilon$ admits one on $r\le k$ labels of the form
\begin{equation}\label{eq:equivariant56}
 \alpha_s\rho(h^{-1})d_j,\qquad h\in H,
 \qquad \rho:H\longrightarrow S_r\text{ a continuous homomorphism},
\end{equation}
with the same uniform error. Here permutations are represented by row-action matrices. In particular the command row is $\rho(u_a^{-1})$, all physical group relations hold on the new register, and initialization may remain random.
\end{theorem}
\begin{proof}
Apply Theorem~\ref{thm:purification55}, whose proof uses only finite convex combinations of decoder values. Denote the resulting initialization by $\beta_s$, its decoder by $D_j$, and its compact joint group by $\mathcal L\subset H\times S_k$, with the product topology and discrete permutation factor. Unused labels may be adjoined when necessary. Let $P$ be the second projection and
\[
                 N=\{\pi\in P:(1,\pi)\in\mathcal L\}.
\]
Then $N$ is normal in $P$. The averaging matrix $Q=|N|^{-1}\sum_{n\in N}n$ is stochastic, satisfies $Q^2=Q$, and commutes with $P$. For $(h,\pi)\in\mathcal L$, every $(h,n\pi)$ belongs to $\mathcal L$. Averaging their error inequalities and using convexity of the norm ball gives
\begin{equation}\label{eq:kernelaverage56}
                   \|\beta_sQ\pi D_j-F_{s,j}(h^{-1})\|_j\le\varepsilon.
\end{equation}

Let $O_1,\ldots,O_r$ be the $N$-orbits in the label set, and let $v_i$ be the uniform probability row on $O_i$. The rows of $Q$ are precisely these uniform orbit rows. Write $V$ for the matrix with rows $v_i$. Every $\pi\in P$ maps $N$-orbits bijectively onto $N$-orbits and therefore induces a permutation $\bar\pi$ with
\begin{equation}\label{eq:orbitintertwining56}
                             V\pi=\bar\pi V.
\end{equation}
The induced action of each $n\in N$ is trivial. Two elements of $\mathcal L$ with the same first coordinate differ by an element of $\{1\}\times N$. Hence $(h,\pi)\mapsto\bar\pi$ descends to a homomorphism $\rho:H\to S_r$. It is continuous: the surjection from the compact group $\mathcal L$ onto the Hausdorff group $H$ is a quotient map, and the original orbit action is continuous into a finite discrete group.

Set $\alpha_s(i)=\sum_{x\in O_i}\beta_s(x)$ and $d_j(i)=v_iD_j$. Then $\alpha_sV=\beta_sQ$ and $d_j=VD_j$. Equation~\eqref{eq:orbitintertwining56} gives
\[
                 \alpha_s\rho(h)d_j=\beta_sQ\pi D_j
                 \quad((h,\pi)\in\mathcal L).
\]
Replace $h$ by $h^{-1}$ in \eqref{eq:kernelaverage56}. This proves \eqref{eq:equivariant56}; all decoder values are legal by convexity, and $r\le k$.
\end{proof}

\begin{corollary}[Intrinsic stationary minimum]\label{cor:intrinsicminimum56}
The number $M_\varepsilon$ is the least cardinality of a finite continuous $H$-set for which some probability initialization for each seed and legal decoder for each query approximate the interface within $\varepsilon$. It is independent of the chosen finite dense alphabet. It is unchanged by replacing every target with
\[
                       h\longmapsto F_{s,j}(\varphi(h)q),
\]
where $\varphi$ is a continuous group automorphism and $q\in H$.
\end{corollary}
\begin{proof}
Theorem~\ref{thm:equivariant56} proves necessity, and restriction to command words proves sufficiency. For the last assertion, replace $\rho$ by $\rho\circ\varphi$ and $\alpha_s$ by $\alpha_s\rho(q^{-1})$ in \eqref{eq:equivariant56}. Applying the inverse change of variables proves equality of the minima.
\end{proof}
The finite quotient in Theorem~\ref{thm:finitequotient55} has a more precise description than its factorial bound. With $\mathcal L,P,N$ as above and $U=\{h:(h,I)\in\mathcal L\}$, the correspondence $(h,\pi)\mapsto(hU,\pi N)$ induces $H/U\simeq P/N$. Thus $[H:U]=|P/N|\le k!$. The orbit action in Theorem~\ref{thm:equivariant56} can have an additional kernel; no equality between this index and the minimal label count is asserted.

\begin{theorem}[Complete clocked finite-quotient classification]\label{thm:completeboundary56}
Assume eventual approximate returns on the compact group $H$. At every $\varepsilon\ge0$, including a positive critical radius, the following are equivalent:
\begin{enumerate}
\item $\sup_NW_{N,\varepsilon}<\infty$;
\item $M_\varepsilon<\infty$;
\item $r(U)\le\varepsilon$ for an open normal subgroup $U$ of $H$.
\end{enumerate}
For each $k<M_\varepsilon$ the all-profile occupation bound \eqref{eq:occupation56} holds, and $W_{N,\varepsilon}$ converges to $M_\varepsilon$. Put
\[
 R_0=\max_{s,j,h\in H}\rad_{Y_j}(F_{s,j}(hH^\circ)).
\]
Then $R_0=\inf_{U\lhd H\text{ open}}r(U)$. Above $R_0$ bounded realizations exist; below it every fixed width has bounded occupation. At $\varepsilon=R_0$ bounded clocked realization exists exactly when the infimum is attained by an open normal subgroup.
\end{theorem}
\begin{proof}
The first two conditions are equivalent by Theorem~\ref{thm:stationarization56}. The second and third are equivalent by Theorem~\ref{thm:finitequotient55}. That proof, including constrained center attainment, uses compactness and convexity of $Y_j$, not finite dimensionality.

For completeness, every open subgroup contains $H^\circ$: the image of the connected identity component in a finite discrete quotient is a singleton. Thus $r(U)\ge R_0$. The quotient $H/H^\circ$ is profinite, so its open normal subgroups form a neighborhood basis at the identity. The inverse images $U$ have intersection $H^\circ$. Given a neighborhood $O$ of $H^\circ$ in $H$, compactness supplies one such $U\subset O$: otherwise their intersections with the closed set $H\setminus O$ have the finite intersection property. Uniform continuity of the finite family of targets gives, for every $\eta>0$, a neighborhood $V$ of $1$ such that $\|F_{s,j}(hkv)-F_{s,j}(hk)\|_j<\eta$ for all $h\in H$, $k\in H^\circ$, and $v\in V$. Choose $U\subset H^\circ V$. Every $F_{s,j}(hU)$ is then within $\eta$ of $F_{s,j}(hH^\circ)$, so $r(U)\le R_0+\eta$. This proves the infimum formula and all remaining assertions.
\end{proof}
This result includes the formerly separate zero and positive boundary cases. The finite-rank proof of Theorem~\ref{thm:zero55} is retained below because it identifies an additional exact translate-space obstruction. No passage from an error strictly larger than $R_0$ to equality is used in Theorem~\ref{thm:completeboundary56}.

\subsection{Length phases without a width overhead}
Let $p\ge1$ be an integer and fix a command with product $a$. Let $J$ be the closure of all products of lengths divisible by $p$. The group argument of Lemma~\ref{lem:phasegroup55}, which uses only the block length, shows that $J$ is open normal, $H/J$ is cyclic of order dividing $p$, and all letters belong to $Ja$. Regard $\mathcal A^p$ as a finite alphabet on $J$. Assume that this block alphabet has eventual approximate returns. In particular this holds when the original nonempty exact-return language has greatest common divisor $p$.

For $0\le r<p$ let $M_{r,\varepsilon}$ be the intrinsic stationary minimum on $J$ for
\[
                         f_{r,s,j}(h)=F_{s,j}(ha^r).
\]
\begin{theorem}[Exact limiting width in each phase]\label{thm:phaseminimum56}
Under the preceding assumptions, for every $r$ and $\varepsilon\ge0$,
\begin{equation}\label{eq:phaseminimum56}
 W_{N,\varepsilon}\le M_{r,\varepsilon}\quad(N\equiv r\pmod p),
 \qquad
 \lim_{\substack{N\to\infty\\N\equiv r\ ({\rm mod}\ p)}}W_{N,\varepsilon}
          =M_{r,\varepsilon}.
\end{equation}
Finite limits are eventually attained. If $k<M_{r,\varepsilon}$, the number of positive cuts of width at most $k$ is bounded uniformly over that horizon phase, with no restriction on other widths. Boundedness in the phase is equivalent to the existence of an open normal subgroup $U\le J$ such that
\begin{equation}\label{eq:phasequotient56}
          \max_{s,j,h\in J}\rad_{Y_j}(F_{s,j}(hUa^r))\le\varepsilon.
\end{equation}
\end{theorem}
\begin{proof}
Suppose first that $M_{r,\varepsilon}=K<\infty$ and choose a physical-equivariant stationary realization of $f_r$ on $K$ labels, with representation $\rho:J\to S_K$. For a fixed $N\equiv r\pmod p$, let $g_t$ be the physical product after $t$ commands and set
\begin{equation}\label{eq:gauge56}
 z_t=a^{N-t}g_ta^{-r}\in J,
 \qquad c_{t,b}=a^{N-t}u_ba^{-(N-t+1)}\in J.
\end{equation}
Normality of $J$ and $u_ba^{-1}\in J$ justify the membership assertions. We have $z_0=a^{N-r}$, $z_t=c_{t,b}z_{t-1}$ when the $t$th letter is $b$, and $z_N=g_Na^{-r}$. Initialize with $\alpha_s\rho(z_0^{-1})$, use the clock-dependent permutation $\rho(c_{t,b}^{-1})$, and keep decoder $d_j$. The row product is $\alpha_s\rho(z_N^{-1})$, so its error from $F_{s,j}(g_N)=f_{r,s,j}(z_N)$ is at most $\varepsilon$. This uses the free external clock but no phase register or partial-block storage. It proves the upper bound with exactly $K$ labels.

Fix $k<M_{r,\varepsilon}$ and a cut residue $0\le\ell<p$. Put $b=(r-\ell)\bmod p$, chosen in $\{0,\ldots,p-1\}$. Restrict the input to the fixed prefix $a^\ell$, arbitrary $p$-letter blocks, and fixed suffix $a^b$. Integrate the prefix into initialization and the suffix into the decoder. The block target on $J$ is
\[
 G_{\ell,s,j}(h)=F_{s,j}(a^b h a^\ell)
     =f_{r,s,j}\bigl(\operatorname{Ad}_{a^b}(h)a^{\ell+b-r}\bigr).
\]
Here $a^{\ell+b-r}\in J$ and conjugation by $a^b$ is an automorphism of $J$. Corollary~\ref{cor:intrinsicminimum56} therefore identifies its stationary minimum with $M_{r,\varepsilon}$. Theorem~\ref{thm:stationarization56} on the block alphabet bounds its positive narrow cuts by a constant $L_\ell$. These cuts are precisely the original cuts in residue $\ell$ between the prefix and suffix, apart from the initial block cut. At most  $\ell+b+1\le2p$ further cuts need be charged for this residue, including small horizons shorter than the chosen prefix and suffix. Summing gives the uniform bound $\sum_{\ell=0}^{p-1}L_\ell+2p^2$. As before, a peak at most $k$ would force $N$ below this bound. This proves the limit and occupation assertions. Finally apply Theorem~\ref{thm:finitequotient55} to $f_r$ on $J$ to obtain \eqref{eq:phasequotient56}.
\end{proof}
If the order of $H/J$ is a proper divisor $q$ of $p$, the minima and quotient-radius criteria repeat with period $q$: multiplication by $a^q\in J$ is a right translation of the target. The earlier phase radii have the same repetition. Neither the physical quotient order nor the return-length greatest common divisor is substituted silently for the other.

\subsection{Compact outputs and finite hidden state}
Theorems~\ref{thm:stationarization56}, \ref{thm:equivariant56}, \ref{thm:completeboundary56}, and \ref{thm:phaseminimum56} hold for norm-compact convex Banach-valued outputs. A single-machine purification even needs only convexity: its new decoder values lie in the convex hull of finitely many old values. Compact output sets are used instead in finite-witness extraction, optimal-error attainment, and constrained-center existence. The representative-function and algebraic sections below retain their explicitly finite-dimensional hypotheses.

Finiteness on the hidden side has a different role. On countably many labels the full simplex is not compact in total variation; mass can escape to infinity, matrix ranks need not admit a useful finite minimum, and a recurrent simplex may have infinitely many vertices. On a general compact convex hidden state space affine automorphisms need not permute a finite vertex set: rotation of the disk already acts continuously on infinitely many extreme points. Thus these proofs give neither a finite-label purification theorem for arbitrary operator kernels nor an online observation theorem. Their terminal-output assertion is strictly the one stated in the models.
